\chapter{Inflation Derivatives}\label{ch:rc:inflation-derivatives}

In October 2022 the UK retail price index stood 14.2\% above its level a year
earlier. Many British pension schemes do not promise their members full
inflation: their increases are the index's annual change floored at zero and
capped at 5\% (2.5\% for rights earned since 2005). Schemes that had hedged those
promises with plain index swaps held hedges that paid the whole 14.2\% on
liabilities that would rise by 5; the difference between an index and a capped,
annually reset index had become the largest number in their risk report. Every
inflation product beyond the zero-coupon swap of One Quant Book 2, chapter 11,
depends on how the index moves from year to year, and that is a question for a
model. This chapter builds the standard one on an analogy with foreign exchange,
computes the convexity that separates year-on-year from zero-coupon products,
adds seasonality, and prices \omterm{def:rc:inflation-derivatives:lpi}{limited price indexation} by simulation.

\section{The foreign-currency analogy}

Treat nominal money as one currency and ``real'' money, a unit of the basket,
as another. A nominal bond pays currency; an index-linked bond pays a basket, worth
$I(T)$ units of currency. The index is the exchange rate between the two, real
discount factors $P_R(t,T)$ are the foreign currency's discount factors, and the
forward index is a forward exchange rate:
\[
I(t,T) = I(t)\,\frac{P_R(t,T)}{P_N(t,T)},
\]
a martingale under the nominal $T$-forward measure. A zero-coupon inflation swap
fixes $(1+k_T)^T = I(0,T)/I(0)$, so today's real curve is read from the swaps:
$P_R(0,T) = P_N(0,T)(1+k_T)^T$.

\begin{definition}[Jarrow--Yildirim model]\label{def:rc:inflation-derivatives:jy}
The \emph{Jarrow--Yildirim model}\index{Jarrow--Yildirim model} (2003) models the
nominal short rate and the real short rate as two Hull--White processes and the
index as a lognormal exchange rate between them:
$dI/I = (n_t-r^R_t)\,dt+\sigma_I\,dW_I$ under the nominal measure, with correlations
between the three drivers. As with a foreign short rate seen from home (the quanto
adjustment of One Quant Book 5, chapter 17), the real rate acquires the drift
$-\rho\,\sigma_R\sigma_I$ under the nominal measure, $\rho$ the correlation of the real
rate with the index.
\end{definition}

The chapter keeps nominal rates deterministic, which isolates the effects of the
real rate and the index; the full model adds terms in the nominal volatility and
its correlations of the same form.

\begin{omfigure}
\begin{tikzpicture}[font=\small,
  box/.style={draw,thick,rounded corners=2pt,align=center,minimum height=1cm,minimum width=3.4cm,inner sep=4pt},
  arr/.style={-{Stealth[length=2.5mm]},thick}]
\node[box,draw=omDef,fill=omDef!8] (n) at (0,0) {nominal economy\\rates $n_t$, bonds $P_N$};
\node[box,draw=omThm,fill=omThm!8] (r) at (8,0) {real economy\\rates $r^R_t$, bonds $P_R$};
\draw[arr,omProp] (r.160) -- node[above=4pt,align=center] {index $I(t)$: currency per basket\\(the exchange rate)} (n.20);
\node[align=center,font=\footnotesize] at (4,-1.35) {domestic currency $\leftrightarrow$ nominal;\quad foreign currency $\leftrightarrow$ real;\quad FX forward $\leftrightarrow$ forward index};
\end{tikzpicture}

\omcaption{The foreign-currency analogy behind inflation models: real money is a
foreign currency whose exchange rate is the price index, so the machinery of
cross-currency pricing, including quanto drifts, carries over.}\label{fig:rc:inflation-derivatives:analogy}
\end{omfigure}

\section{Zero-coupon versus year-on-year}

\begin{definition}[Year-on-year inflation swap]\label{def:rc:inflation-derivatives:yoy}
A \emph{year-on-year inflation swap}\index{year-on-year inflation swap} exchanges,
each year $i$, the index's annual change $I(T_i)/I(T_{i-1})-1$ against a fixed
rate, on a notional that does not grow with the index.
\end{definition}

A zero-coupon swap pays one ratio $I(T)/I(0)$ at $T$: its value needs only the forward
index. A year-on-year leg pays a ratio of two future index values, paid at the later
date; its value needs the expectation of the ratio under the nominal measure of the
payment date, and the ratio's numerator and denominator are correlated through the
real rate.

\begin{definition}[Year-on-year convexity adjustment]\label{def:rc:inflation-derivatives:yoyadj}
The \emph{year-on-year convexity adjustment}\index{year-on-year convexity adjustment}
of period $i$ is the difference between $\E^{T_i}_N[I(T_i)/I(T_{i-1})]$ and the ratio of
forward indices $I(0,T_i)/I(0,T_{i-1})$.
\end{definition}

\begin{proposition}[Year-on-year expectation]\label{prop:rc:inflation-derivatives:yoy}
In the \omterm{def:rc:inflation-derivatives:jy}{Jarrow--Yildirim model} with deterministic nominal rates, real rate
$r^R = f_R(0,t)+x_t$ with Hull--White parameters $(\kappa,\sigma_R)$ and correlation $\rho$
between the real rate and the index,
\[
\E_N\Bigl[\frac{I(S)}{I(T)}\Bigr] = \frac{I(0,S)}{I(0,T)}\,e^{C(T,S)},\quad
C(T,S) = -B(T,S)\Bigl(\frac{\sigma_R^2}{2\kappa^2}\bigl(1-e^{-\kappa T}\bigr)^2-\rho\,\sigma_R\sigma_I\,B(0,T)\Bigr).
\]
\end{proposition}
\begin{proof}
Conditioning at $T$ gives $\E_T[I(S)] = I(T)\,P_R(T,S)/P_N(T,S)$, so the
expectation of the ratio is $\E_N[P_R(T,S)]/P_N(T,S)$. With
\[
P_R(T,S) = \frac{P_R(0,S)}{P_R(0,T)}\,e^{-Bx_T-\frac12B^2y(T)}
\]
and $x_T$ normal with variance $y(T)$ and nominal-measure mean $\mu(T)-\rho\sigma_R\sigma_IB(0,T)$,
where $\mu(T) = \int_0^Te^{-\kappa(T-u)}y(u)\,du = \frac{\sigma_R^2}{2\kappa^2}(1-e^{-\kappa T})^2$, take
$\E[e^{-Bx_T}]$.
\end{proof}

\omcode{code/firm/inflopt/firm_inflopt.py}{63}{70}{The year-on-year expectation and its convexity adjustment.}\label{lst:rc:inflation-derivatives:yoy}

\begin{example}[Sterling year-on-year convexity]\label{ex:rc:inflation-derivatives:conv}
On illustrative sterling curves (SONIA zero rates of 3.6--4.2\%, RPI zero-coupon
swaps of 3.2--3.4\%) with $\kappa=5\%$, $\sigma_R=80$ basis points, $\sigma_I=1.5\%$ and
$\rho=0.2$, the adjustment is zero for the first year, $-1.8$ basis points for the fourth,
$-15.2$ for the tenth (forward 3.42\%, year-on-year expectation 3.27\%) and $-42.5$ for the
nineteenth (\cref{fig:rc:inflation-derivatives:convexity}). The ten-year year-on-year swap
rate is 3.22\% against a ten-year zero-coupon rate of 3.28\%. A simulation of the model
reproduces the tenth year's expectation, 1.03268, at 1.03277 with a standard error of
0.00013.
\end{example}

\begin{omfigure}
\begin{tikzpicture}
\begin{axis}[width=11.5cm,height=5.4cm,
  xlabel={year $i$ of the annual change},ylabel={rate (\%)},
  tick label style={font=\small},label style={font=\small},
  xmin=0.5,xmax=20.5,ymin=2.8,ymax=3.6,grid=major,grid style={gray!25},
  legend columns=2,legend style={font=\footnotesize,at={(0.5,-0.25)},anchor=north,draw=none,fill=none,/tikz/every even column/.append style={column sep=8pt}},
  table/col sep=comma]
\addplot[omDef,thick,mark=*,mark size=1.4pt] table[x=i,y=fwd]{figdata/rates-credit-risk/11-inflation-derivatives/convexity.csv};
\addplot[omThm,very thick,mark=square*,mark size=1.4pt] table[x=i,y=yoy]{figdata/rates-credit-risk/11-inflation-derivatives/convexity.csv};
\legend{ratio of forward indices,year-on-year expectation}
\end{axis}
\end{tikzpicture}

\omcaption{Annual inflation implied by zero-coupon swaps (ratio of forward indices)
and the model's expectation of each year's change, which a year-on-year swap pays.
The gap, the convexity adjustment, grows with the start of the year. Data:
illustrative sterling curves, Jarrow--Yildirim parameters of
\cref{ex:rc:inflation-derivatives:conv}; the chapter's
tutorial.}\label{fig:rc:inflation-derivatives:convexity}
\end{omfigure}

The sign and size depend on the correlation: at the tenth year the adjustment is
$-20.6$, $-16.4$ and $-12.2$ basis points for $\rho = -0.5$, $0$ and $0.5$. With no real-rate
volatility it vanishes whatever $\rho$; the index volatility alone does not create it.

\section{Seasonality in the model}

Price indices are not seasonally adjusted, and their monthly changes follow the
calendar: sales, energy tariffs, holidays (One Quant Book 2, chapter 11). An
inflation swap or bond references a given month's index, so the forward index must
carry the pattern between annual points while matching them.

\begin{method}[Seasonal forward index]\label{met:rc:inflation-derivatives:season}
Estimate the average excess of each calendar month's log change over its year's
mean (the seasonal factors, summing to zero); build the forward index month by month
as the smooth forward times the cumulated factors. Annual points are unchanged, the
path between them bends, and any product fixing on a month other than the swaps'
reference months moves.
\end{method}

\begin{example}[The US pattern]\label{ex:rc:inflation-derivatives:season}
On US CPI (not seasonally adjusted) over the complete years 2010--2025, the monthly
log change exceeds its year's average by 0.18, 0.26 and 0.31 percentage points in
January, February and March and falls short by 0.36 and 0.35 in November and
December (\cref{fig:rc:inflation-derivatives:seasonal}). Added to a forward index
that is smooth at 3\% a year, the pattern lifts it by 1.0 point by the end of May and
1.1 by the end of June, and brings it back to the smooth path at each year-end.
\end{example}

\begin{omfigure}
\begin{tikzpicture}
\begin{axis}[width=11.5cm,height=5.2cm,
  xlabel={months from today (January)},ylabel={forward index (base 100)},
  tick label style={font=\small},label style={font=\small},
  xmin=0,xmax=24,ymin=99.5,ymax=106.5,xtick={0,3,6,9,12,15,18,21,24},grid=major,grid style={gray!25},
  legend columns=2,legend style={font=\footnotesize,at={(0.5,-0.25)},anchor=north,draw=none,fill=none,/tikz/every even column/.append style={column sep=8pt}},
  table/col sep=comma]
\addplot[omDef,thick,dashed] table[x=m,y=smooth]{figdata/rates-credit-risk/11-inflation-derivatives/seasonal.csv};
\addplot[omThm,very thick,mark=*,mark size=1.2pt] table[x=m,y=seasonal]{figdata/rates-credit-risk/11-inflation-derivatives/seasonal.csv};
\legend{smooth at 3\% a year,with the US seasonal pattern}
\end{axis}
\end{tikzpicture}

\omcaption{A two-year forward index at 3\% a year, smooth and with the monthly
seasonal factors of US CPI (complete years 2010--2025). Both pass through the same
annual points. Source: US Bureau of Labor Statistics via FRED (series CPIAUCNS, One
Quant Book 2's data file); the chapter's
tutorial.}\label{fig:rc:inflation-derivatives:seasonal}
\end{omfigure}

\section{Inflation options and limited price indexation}

\begin{definition}[Inflation cap]\label{def:rc:inflation-derivatives:cap}
An \emph{inflation cap}\index{inflation cap} is a strip of options each paying
$\max(I(T_i)/I(T_{i-1})-1-K,0)$ at $T_i$: a cap on year-on-year inflation, with floors
(the deflation floor of One Quant Book 2, chapter 11, for the zero-coupon case) defined
likewise.
\end{definition}

In the model the annual ratio is lognormal: its log-variance is the index variance
over the year plus the variance of the year's real-rate integral, less twice their
covariance, and the caplet is Black's formula on the year-on-year expectation.

\begin{example}[A ten-year cap and floor]\label{ex:rc:inflation-derivatives:caps}
On the sterling curves, a year-on-year cap at 5\% on years two to ten costs 181.1 basis
points of notional; a floor at 0\% costs 45.8.
\end{example}

\begin{definition}[Limited price indexation]\label{def:rc:inflation-derivatives:lpi}
\emph{Limited price indexation}\index{limited price indexation} (LPI) uprates an amount
each year by the index's annual change floored at $f$ and capped at $c$, compounding
the result: after $N$ years it has grown by $\prod_{i=1}^N\bigl(1+\min(\max(I_i/I_{i-1}-1,f),c)\bigr)$.
UK pension increases follow it by statute, with $[0\%,5\%]$ for rights earned before 6
April 2005 and $[0\%,2.5\%]$ after.
\end{definition}

Compounding capped and floored changes is path-dependent: the payoff is not a sum of
caplets. It is priced by simulating the year-on-year ratios jointly.

\omcode{code/firm/inflopt/firm_inflopt.py}{119}{127}{The limited-price-indexation leg by simulation of annual index ratios.}\label{lst:rc:inflation-derivatives:lpi}

\begin{omfigure}
\begin{tikzpicture}
\begin{axis}[width=11.5cm,height=5.2cm,
  xlabel={growth factor after twenty years},ylabel={share of paths (\%)},
  tick label style={font=\small},label style={font=\small},
  xmin=1.0,xmax=3.1,ymin=0,ymax=5.5,xtick={1,1.5,2,2.5,3},grid=major,grid style={gray!25},
  legend columns=2,legend style={font=\footnotesize,at={(0.5,-0.3)},anchor=north,draw=none,fill=none,/tikz/every even column/.append style={column sep=8pt}},
  table/col sep=comma]
\addplot[omDef,thick,const plot mark mid] table[x=x,y=full]{figdata/rates-credit-risk/11-inflation-derivatives/lpihist.csv};
\addplot[omThm,very thick,const plot mark mid] table[x=x,y=lpi]{figdata/rates-credit-risk/11-inflation-derivatives/lpihist.csv};
\legend{index ratio $I(20)/I(0)$,{LPI $[0\%,5\%]$}}
\end{axis}
\end{tikzpicture}

\omcaption{Distribution of the twenty-year growth of an uncapped index and of its
$[0\%,5\%]$ limited-price-indexed version, 40\,000 simulated paths of the sterling model:
the cap removes the right tail, the floor the rare deflations. Data: the chapter's
tutorial.}\label{fig:rc:inflation-derivatives:lpi}
\end{omfigure}

\begin{dated}{2026-09}{UK inflation in 2022 and after}\label{dat:rc:inflation-derivatives:uk}
The RPI's annual change was 14.2\% in October 2022 (CPI 11.1\%, CPIH 9.6\%). The
government and the UK Statistics Authority decided in November 2020 that the RPI will
be brought into line with CPIH methods and data from February 2030, after the last
index-linked gilt that would otherwise be affected, which moves the RPI swap curve
beyond 2030. Statutory pension increases remain limited to 5\% (pre-2005 rights) or
2.5\%.
\end{dated}

\section{Tutorial: year-on-year products}\label{tut:rc:inflation-derivatives:tutorial}

\begin{tutorial}
\textbf{Goal.} Build nominal and real curves, compute year-on-year convexity, price
caps and an LPI leg, and put seasonality into a forward index. \textbf{End state:}
\cref{fig:rc:inflation-derivatives:convexity,fig:rc:inflation-derivatives:seasonal,fig:rc:inflation-derivatives:lpi}.
\begin{tutsteps}
\item \textbf{Curves}: \texttt{real\_curve(nominal, years, zc\_rates)} from the RPI
zero-coupon swaps.
\item \textbf{Convexity}: \texttt{convexity\_table()}, \texttt{rho\_effect()} and
\texttt{mc\_check()}.
\item \textbf{Options}: \texttt{caps()} and \texttt{lpi()}.
\item \textbf{Seasonality}: \texttt{us\_seasonals()} with Book~2's
\texttt{seasonal\_factors}; \texttt{fig\_rc\_inflation.py} writes the charts.
\end{tutsteps}
\textbf{What to change next.} Price the LPI leg with the cap at 2.5\% and a floor of
$-100\%$ (no floor) and split the value between the cap and the floor; make the nominal
rate stochastic in the simulation and measure the change in the year-on-year swap
rate.
\end{tutorial}

\section{Build: inflation options}\label{bld:rc:inflation-derivatives:inflopt}

\begin{build}
\bfield{Purpose} The inflation book of the miniature firm: year-on-year swaps, caps
and floors, and the LPI legs of pension-scheme hedges.
\bfield{Interface} \texttt{LogLinearCurve}, \texttt{real\_curve};
\texttt{JYDet(nominal, real, kappa, sigma\_r, sigma\_i, rho)} with \texttt{yoy\_ratio},
\texttt{convexity}, \texttt{yoy\_log\_variance}, \texttt{yoy\_caplet},
\texttt{yoy\_swap\_rate}, \texttt{zc\_swap\_rate}, \texttt{simulate\_ratios},
\texttt{lpi\_leg}.
\bfield{Rules} Deterministic nominal rates; the real curve reprices the zero-coupon
swaps exactly; simulation under the nominal measure with the quanto drift; Book~2's
\texttt{firm\_breakeven} for seasonal factors.
\bfield{Acceptance tests} \texttt{code/firm/inflopt/tests/}: the model reprices the
zero-coupon swap rate and has no adjustment in the first year or without real-rate
volatility; the analytic year-on-year expectation matches simulation; cap minus floor
is the forward; an LPI leg with infinite bounds is the index, and the simulated index
matches the real discount factor.
\bfield{Stretch} Stochastic nominal rates (the full model); a smile for year-on-year
caplets; seasonality in the model's forward index at monthly fixings.
\end{build}

\begin{omsources}
\item R.\ Jarrow and Y.\ Yildirim, ``Pricing Treasury Inflation Protected Securities
and related derivatives using an HJM model'', \emph{Journal of Financial and
Quantitative Analysis} 38(2), 2003.
\item F.\ Mercurio, ``Pricing inflation-indexed derivatives'', \emph{Quantitative
Finance} 5(3), 2005.
\item ONS, \emph{Consumer price inflation, UK: October 2022}; HM Treasury and UK
Statistics Authority, response on RPI reform, November 2020; Pensions Act 1995,
section 51.
\end{omsources}

\section{Exercises}

\begin{exercise}[$\star$]\label{exo:rc:inflation-derivatives:1}
A ten-year zero-coupon RPI swap is at 3.28\% and the ten-year nominal discount factor
is 0.6771. Give the ten-year real discount factor and the forward index ratio.
\end{exercise}

\begin{exercise}[$\star$]\label{exo:rc:inflation-derivatives:2}
Why is there no convexity adjustment for the first year of a year-on-year swap?
\end{exercise}

\begin{exercise}[$\star$]\label{exo:rc:inflation-derivatives:3}
A pension increase with LPI $[0\%,5\%]$ follows an index that rose 14.2\%. By how much
does the pension rise? And with $[0\%,2.5\%]$?
\end{exercise}

\begin{exercise}[$\star\star$]\label{exo:rc:inflation-derivatives:4}
Explain the sign of the adjustment in \cref{ex:rc:inflation-derivatives:conv} and why it
becomes less negative as $\rho$ rises.
\end{exercise}

\begin{exercise}[$\star\star$]\label{exo:rc:inflation-derivatives:5}
Read \cref{fig:rc:inflation-derivatives:seasonal}: a swap references the May index and
another the December index of next year. Which has the higher forward, relative to the
smooth path, and why?
\end{exercise}

\begin{exercise}[$\star\star$]\label{exo:rc:inflation-derivatives:6}
Why is an LPI leg not a strip of year-on-year caps and floors?
\end{exercise}

\begin{exercise}[$\star\star\star$]\label{exo:rc:inflation-derivatives:7}
\emph{Coding.} With common random numbers, measure the change in the twenty-year LPI
$[0\%,5\%]$ leg when $\sigma_I$ rises from 1.5\% to 2.5\%, and when $\sigma_R$ rises from 0.80\% to
0.90\%.
\end{exercise}

\begin{exercise}[$\star\star\star$]\label{exo:rc:inflation-derivatives:8}
\emph{Find the flaw.} ``The year-on-year swap rate equals the zero-coupon swap rate,
since both average the same forward inflation.''
\end{exercise}

\section{Problem: The Pension Fund's LPI Swap}

\begin{problem}[{Weekend problem --- hedging capped pension increases}]\label{pb:rc:inflation-derivatives:1}
A pension scheme's liabilities rise each year by RPI with LPI $[0\%,5\%]$. Its adviser
proposes a twenty-year swap that pays the LPI growth of a notional at year twenty, against
the uncapped index. Use the sterling curves and parameters of the chapter, 40\,000 paths.

\medskip\noindent\textbf{Part I --- The legs.}
\begin{enumerate}
\item Give today's value of the uncapped index leg, per unit, and its exact value.
\item Give the value of the LPI $[0\%,5\%]$ leg.
\item Give the value of an LPI $[0\%,2.5\%]$ leg.
\item Why is the LPI leg cheaper, and which of its two options dominates?
\item What does the difference mean for a scheme hedged with uncapped swaps?
\end{enumerate}

\medskip\noindent\textbf{Part II --- Sensitivities.}
\begin{enumerate}[resume]
\item Give the change in the LPI leg when $\sigma_I$ rises by one point.
\item Give the change when $\sigma_R$ rises by ten basis points.
\item Is the scheme, after the swap, long or short volatility?
\item Why are the sensitivities computed on common random numbers?
\item Which market instruments would hedge them?
\end{enumerate}

\medskip\noindent\textbf{Part III --- 2022.}
\begin{enumerate}[resume]
\item With the index up 14.2\%, what increase did the $[0\%,5\%]$ liabilities receive?
\item What did an uncapped hedge pay on the same year?
\item Why did 2022 turn the scheme's hedge into a speculative position?
\item What does the RPI reform of 2030 change for the LPI leg?
\item How would the problem change for an index with a floor of $-100\%$?
\end{enumerate}

\medskip\noindent\textbf{Part IV --- Judgement.}
\begin{enumerate}[resume]
\item Who can take the other side of LPI swaps?
\item Why is the model's correlation $\rho$ a risk the scheme cannot see in quotes?
\item What would you report to the trustees?
\item State the \emph{named result}: the values of the uncapped and LPI legs, and the
LPI leg's index-volatility sensitivity.
\item In one sentence: what does a cap on annual changes do to a twenty-year promise?
\end{enumerate}
\end{problem}

\section{Interview questions}

\begin{interviewq}[$\star$ \iqroles{researcher, bank}]\label{iq:rc:inflation-derivatives:1}
Explain the foreign-currency analogy for inflation.
\end{interviewq}

\begin{interviewq}[$\star\star$ \iqroles{researcher}]\label{iq:rc:inflation-derivatives:2}
Why is a year-on-year swap not priced off the zero-coupon curve alone?
\end{interviewq}

\begin{interviewq}[$\star\star$ \iqroles{trader}]\label{iq:rc:inflation-derivatives:3}
How does seasonality enter an inflation swap's price, and where do you get it?
\end{interviewq}

\begin{interviewq}[$\star\star$ \iqroles{trader, bank}]\label{iq:rc:inflation-derivatives:4}
What is \omterm{def:rc:inflation-derivatives:lpi}{limited price indexation}, and why does it make pension hedging hard?
\end{interviewq}

\begin{interviewq}[$\star\star\star$ \iqroles{researcher, risk}]\label{iq:rc:inflation-derivatives:5}
Which parameters of the \omterm{def:rc:inflation-derivatives:jy}{Jarrow--Yildirim model} are observable in the market, and which
are not?
\end{interviewq}

\begin{interviewq}[$\star\star\star$ \iqroles{developer}]\label{iq:rc:inflation-derivatives:6}
Your simulated index does not reprice the zero-coupon swaps. What do you check?
\end{interviewq}
